\begin{itemize}
    \item 任意圖
        \begin{itemize}
            \item 最大匹配 + 最小邊覆蓋 = $n$ （不能有孤點）
            \item  點覆蓋的補集是獨立集。最小點覆蓋 + 最大獨立集 = $n$ 
            \item  $w($最小權點覆蓋$)$ + $w($最大權獨立集$) = \sum{w_v}$
            \item  （帶點權的二分圖可以用最小割解，構造請參考 Dinic.cpp 的 construct ）
        \end{itemize}
    \item 二分圖
        \begin{itemize}
            \item 最小點覆蓋 = 最大匹配 = $n$ - 最大獨立集
        \end{itemize}
    \item 只有邊帶權的二分圖
        \begin{itemize}
            \item w-vertex-cover（帶權點覆蓋）：每條邊的兩個連接點被選中的次數總和至少要是 $w_e$。
            \item w-weight matching（帶權匹配） 
            \item minimum vertex count of w-vertex-cover = maximum weight count of w-weight matching（一個點可以被選很多次，但邊不行）
        \end{itemize}
    \item 點、邊都帶權的二分圖的定理
        \begin{itemize}
            \item b-matching：每條邊 $e$ 可以被選 $x_e$ 次（非負整數），每個點 $v$ 都要滿足 $\sum_{e \ni v}{x_e}\leq b_v$ ，權重為 $\sum{w_e x_e}$ 。
            \item The maximum w-weight of a b-matching equals the minimum b-weight of vertices in a w-vertex-cover.
        \end{itemize}
\end{itemize}
